% !TeX root = ../thesis.tex
\chapter{电连接器机器人自主装配系统集成与综合实验研究}

\section{引言}
本章针对电连接器机器人自主抓取与装配系统的工程实现问题，开展系统集成与综合实验研究。首先搭建缩比航空机舱实验平台，完成机械臂、视觉传感器、力/力矩传感器和连接器工装的硬件集成，并建立相机、末端执行器与机器人基坐标系之间的标定关系；随后开展感知---规划子系统联合实验，验证输出的目标身份、精位姿和 \textit{CableMask} 对抓取规划的支撑作用；进一步开展装配过程力控实验，分析在柔性拖缆干扰下的接触响应特性；最后通过连续作业实验，对系统在作业节拍、装配一致性、成功率和人工干预次数等方面的综合性能进行评估。

\section{实验平台搭建与系统集成}
系统级实验首先依赖稳定可靠的平台条件。若硬件布局不统一、坐标关系不准确或软件接口不稳定，感知结果就无法可靠传递给规划与控制模块，后续实验也难以具有重复性和可比性。狭窄空间、复杂遮挡和柔性拖缆并存的任务特点要求实验平台尽可能贴近真实航空机舱作业条件。本节围绕实验平台构建、系统标定和软件集成展开研究，建立硬件环境、坐标变换关系和 ROS 数据流。

\subsection{航空机舱模拟环境与硬件构成}
航空制造现场的狭窄舱段对机器人系统提出“可观测、可到达、可测量且可重复”的工程约束。为在有限空间内实现电连接器自主抓取与装配，实验平台以缩比机舱作为物理载体，集成工业机械臂、RGB-D 视觉传感器、六维力/力矩测量模块、末端夹持机构与连接器工装，使几何边界与柔性线缆干扰能够在同一实验链路中被稳定复现。

\begin{figure}[htbp]
  \centering
  \fbox{\parbox[c][0.22\textheight][c]{0.9\textwidth}{\centering 图5-1\\（空白占位图）}}
  \caption{电连接器机器人自主装配实验平台硬件集成结构图}
  \label{fig:ch5:platform-hw}
\end{figure}

如图 \ref{fig:ch5:platform-hw} 所示，硬件布局围绕“位姿求解—抓取规划—装配力控”的连续链路进行组织。UR5e 机器人提供 850\,mm 工作半径与 0.03\,mm 重复定位精度，使末端在狭窄通道中完成接近、夹持与插接动作的可重复初始化；Robotiq 2F-85 与 3D 打印 TPU 柔性指尖的组合通过两指自适应夹持降低壳体局部应力集中风险，从而减小由柔性线缆干扰引发的夹持偏差累积。

\begin{table}[htbp]
  \centering
  \caption{实验平台主要硬件参数配置}
  \label{tab:ch5:hw-params}
  \begin{tabular}{|p{2.2cm}|p{3.0cm}|p{8.9cm}|}
    \hline
    硬件类别 & 型号/规格 & 关键参数 \\
    \hline
    工业机械臂 & UR5e & 工作半径 850\,mm；额定负载 5\,kg；重复定位精度 0.03\,mm \\
    \hline
    末端夹持机构 & Robotiq 2F-85 + 3D 打印 TPU 柔性指尖 & 两指自适应夹持；适配电连接器壳体夹持区域 \\
    \hline
    RGB-D 视觉传感器 & Intel RealSense D435i & Eye-in-Hand 安装；分辨率 1280$\times$720；有效工作距离 0.2\,m$\sim$1.0\,m \\
    \hline
    六维力/力矩传感器 & ATIGamma & $F_x/F_y$ 量程 130\,N；$F_z$ 量程 400\,N；$T_x/T_y/T_z$ 量程 10\,N$\cdot$m；分辨率 0.05\,N；采样频率 1000\,Hz \\
    \hline
    缩比机舱工装 & 定制缩比机舱实验平台 & 工作空间 600\,mm$\times$600\,mm$\times$500\,mm；设置 3 个连接器工位；线缆长度约 1.5\,m \\
    \hline
    连接器样件 & D38999/26WC35PN 系列 & 典型插接尺度约 8\,mm \\
    \hline
    上位计算平台 & Intel Core i7-10700K + 32\,GB RAM + NVIDIA RTX3070 GPU & 支持视觉推理、规划计算与装配控制算法运行 \\
    \hline
  \end{tabular}
\end{table}

\FloatBarrier

表 \ref{tab:ch5:hw-params} 给出了平台关键硬件参数，并反映了传感器精度与时间对齐对闭环实验的支撑作用。ATIGamma 六维力/力矩传感器以 1000\,Hz 采样频率对装配阶段的接触载荷与力矩变化进行在线测量；RGB-D 视觉传感器在 Eye-in-Hand 架构下提供彩色图与深度图，为目标身份识别、精位姿估计以及柔性线缆形态感知提供数据来源。视觉更新与力觉反馈的时间尺度匹配，能够避免插接阶段由于测量延时造成的误差累积，使接触调节过程具有可重复的观测条件。

同时，缩比机舱工作空间尺度（600\,mm$\times$600\,mm$\times$500\,mm）与线缆长度约 1.5\,m 使柔性拖缆在抓取抬升、转运与预装配对准过程中形成具有代表性的时变几何约束。柔性线缆形态由深度数据支持的线缆分割结果（\textit{CableMask}）进行三维映射，为抓取与运动规划阶段提供可用于碰撞约束/避障评估的环境先验，从而保证柔性线缆干扰、规划约束与装配控制反馈之间的接口一致性。

\FloatBarrier

\subsection{系统标定与精度分析}
在系统集成完成后，需要建立视觉测量与机器人控制之间的一致坐标映射关系。为确保视觉位姿能够稳定转化为机械臂末端控制目标，系统采用 \textit{Eye-in-Hand} 结构进行手眼标定，求取相机坐标系与末端坐标系的刚体变换，并对整体映射精度进行量化评估。

以校准目标在坐标系 \{O\} 下的位姿为例，基座系 \{B\} 下位姿可由下式描述：
$$
{}^{B}T_{O} = {}^{B}T_{E}\cdot {}^{E}T_{C}\cdot {}^{C}T_{O},
$$
其中，${}^{B}T_{E}$ 由机器人正运动学获得，${}^{E}T_{C}$ 为相机与末端之间的固定刚体变换，${}^{C}T_{O}$ 由标定靶的视觉观测得到。

\begin{figure}[htbp]
  \centering
  \fbox{\parbox[c][0.22\textheight][c]{0.9\textwidth}{\centering 图5-2\\（空白占位图）}}
  \caption{末端相机、六维力传感器与夹爪安装关系图}
  \label{fig:ch5:eyeinhand-setup}
\end{figure}

\textbf{（1）标定方法过程（Eye-in-Hand + Tsai-Lenz）。}本系统采用 Eye-in-Hand 手眼标定结构。相机固定安装于机械臂末端，使相机与末端之间的刚体变换在多次运动中保持不变。标定过程中，机器人在预设末端位姿集合上运动，采集标定靶图像并得到校准目标的相机系观测；同时由每个位姿对应的正运动学计算基座系与末端之间的变换。Tsai-Lenz 方法基于多组位姿约束解算相机与末端的刚体变换，并利用最小二乘意义降低噪声对标定结果的影响。为保证几何条件与实际装配观测一致，本次标定采样采用 20 组末端位姿，并在相机有效成像范围内完成标定靶采集。

\begin{figure}[htbp]
  \centering
  \fbox{\parbox[c][0.22\textheight][c]{0.9\textwidth}{\centering 图5-3\\（空白占位图）}}
  \caption{系统坐标系关系与手眼标定示意图}
  \label{fig:ch5:handeye}
\end{figure}

\textbf{（2）精度分析与误差量化。}为验证标定结果的可用性，采用重投影一致性作为误差指标，并将像素级误差映射为空间位置误差。表 \ref{tab:ch5-calib-acc} 汇总了当前实验设置下的量化结果：20 组采样下重投影误差为 0.42\,pixels，对应空间位置误差约为 0.35\,mm。该误差水平反映视觉位姿到基坐标控制目标的初始映射精度，可为后续抓取与插接对位姿初始化提供亚毫米量级的稳定起点。

同时，力/力矩反馈在装配阶段承担接触载荷检测任务。系统对传感器零点进行载荷重心补偿与零点漂移校正，校正后零漂水平约为 0.2\,N。该结果说明接触前的力反馈基线已得到有效抑制，能够减少零漂导致的误判触发，从而保证装配控制过程的可靠性与可重复性。

\begin{table}[htbp]
  \centering
  \begingroup
  \renewcommand{\thetable}{5-3}
  \caption{系统标定结果与精度分析结果}
  \endgroup
  \label{tab:ch5-calib-acc}
  \begin{tabular}{|p{4.0cm}|p{3.0cm}|p{7.0cm}|}
    \hline
    指标 & 结果 & 说明 \\
    \hline
    手眼标定方法 & Tsai-Lenz & 20 组采样数据 \\
    \hline
    重投影误差 & 0.42 pixels & 对应空间位置误差约 0.35\,mm \\
    \hline
    力传感器零漂水平 & 约 0.2\,N & 载荷重心补偿与零点漂移校正后 \\
    \hline
  \end{tabular}
\end{table}

\FloatBarrier

\subsection{软件体系与数据流闭环}
\begin{figure}[htbp]
  \centering
  \fbox{\parbox[c][0.22\textheight][c]{0.9\textwidth}{\centering 图5-4\\（空白占位图）}}
  \caption{基于 ROS 的系统软件架构与数据流向图}
  \label{fig:ch5:ros-arch}
\end{figure}

本节的软件集成以 ROS 话题与服务为通信骨架，在缩比机舱的受限空间内建立从视觉观测到装配执行的闭环链路。图 \ref{fig:ch5:ros-arch} 给出了各节点之间的数据传递关系与阶段触发逻辑，其中视觉输出驱动运动生成，装配执行阶段融合力觉反馈完成稳定过渡。

为了量化该闭环在多时间尺度上的协同一致性，表 \ref{tab:ch5:freqs} 给出各传感器与节点的运行频率，并对应识别、规划、执行与复位等阶段的启动与结束条件。视觉链路以 RGB 与深度的更新节奏提供目标位姿与线缆状态刷新，规划端在完成状态切换后生成 joint\_trajectory，执行端依据指令刷新频率保持末端运动连续性；力觉链路在接触阶段以高采样频率提供力误差输入。

\begin{table}[htbp]
  \centering
  \begingroup
  \renewcommand{\thetable}{5-2}
  \caption{传感器与系统运行频率配置}
  \endgroup
  \label{tab:ch5:freqs}
  \begin{tabular}{|p{2.6cm}|p{4.2cm}|p{4.7cm}|}
    \hline
    模块/传感器 & 运行频率/方式 & 对应阶段与作用 \\
    \hline
    RGB & 30\,fps & 感知输入与目标实例识别触发 \\
    \hline
    深度图像 & 30\,fps & 位姿求解所需深度反投影与障碍表示更新 \\
    \hline
    力/力矩 & 1000\,Hz & 装配接触载荷实时反馈与零漂基线校验 \\
    \hline
    vision\_node & 平均 65\,ms/帧（约 15\,Hz） & 输出目标位姿量与线缆状态量进入规划准备 \\
    \hline
    planning\_node & 任务触发 & 生成阶段轨迹 joint\_trajectory 并交付控制执行 \\
    \hline
    control\_node & 500\,Hz & 轨迹执行期间融合视觉前馈与力觉反馈输出速度指令 \\
    \hline
    state\_manager & 事件驱动 & 管理阶段切换与复位触发条件 \\
    \hline
  \end{tabular}
\end{table}

\FloatBarrier
从阶段过程看，识别阶段由 vision\_node 在约 15\,Hz 的输出节奏完成 connector\_pose 与 cable\_vector 的更新；当状态管理节点触发进入规划生成时，planning\_node 使用最新的目标位姿与障碍表示生成 joint\_trajectory，并将其交付给 control\_node。控制执行阶段由 control\_node 以 500\,Hz 发布速度指令，在单次视觉更新间隔内形成多次控制校正，使轨迹跟踪误差不会因观测延迟而累积。进入插接/接触阶段后，力/力矩传感器以 1000\,Hz 采样更新接触载荷，驱动力位混合调节，使接触非线性出现时仍能维持阻抗控制的收敛性。

在缩比机舱与柔性线缆干扰条件下，插接接触过程呈现明显的时变非线性，若感知更新、规划生成与控制发布之间缺乏匹配的时间尺度，误差会在接口处累积并放大为接触力估计滞后与末端轨迹抖动。表 \ref{tab:ch5:freqs} 给出的多速率通信与节点触发机制，使得视觉观测输出能够按固定节奏进入规划准备，规划结果在控制执行阶段保持一致的目标更新边界，同时力觉反馈以高采样率驱动阻抗调节。该时序耦合为后续系统级综合验证提供了可复现的闭环接口条件，从而保证不同工况下接触调节过程的可比性与稳定性。

\section{感知---规划子系统联合验证实验}
为验证狭窄机舱场景下“感知输出—规划输入”的系统接口有效性，联合实验以电连接器抓取任务为载体，将任务难点集中在目标身份混淆、目标位姿受遮挡影响以及尾部柔性线缆形成时变障碍三类问题。实验设计强调分段推进：首先构建遮挡工况以覆盖不同可见性条件，其次完成线缆先验从像素表征到三维障碍表征的映射，再对比基线方法与本文方法在同一接口设置下的联合效果，并用定量指标与轨迹/耗时/失效模式的联合解释来完成论证闭环。

\begin{figure}[htbp]
  \centering
  \fbox{\parbox[c][0.22\textheight][c]{0.9\textwidth}{\centering 图5-5\\（空白占位图）}}
  \caption{轻度、中度与重度遮挡典型实验场景图}
  \label{fig:ch5:occ-scenes}
\end{figure}

如图 \ref{fig:ch5:occ-scenes} 所示，实验分别构建轻度、中度与重度遮挡三类典型工况。遮挡等级的变化同时影响目标关键区域可见性与线缆在局部空间中的分布密度，因此可在同一任务评价框架下检验联合系统对“观测不确定性增加”和“环境可通行域收缩”的鲁棒性。

\begin{figure}[htbp]
  \centering
  \fbox{\parbox[c][0.22\textheight][c]{0.9\textwidth}{\centering 图5-6\\（空白占位图）}}
  \caption{CableMask 到 Octomap 障碍地图的构建过程图}
  \label{fig:ch5:cablemask-octomap}
\end{figure}

如图 \ref{fig:ch5:cablemask-octomap} 所示，线缆先验由 CableMask 表征得到三维障碍约束，具体包括深度映射得到相机系点云、点云滤波抑制噪声点与离群点、体素化形成离散占据，再在规划地图结构中以 Octomap 的占据概率形式供碰撞检测调用。该接口映射将线缆由“视觉平面信息”转化为“可验证的碰撞约束”，使规划搜索能够在相同的几何语义下进行。

在对比实验中，Baseline 方法仅采用相对有限的感知输入，而本文方法在同一遮挡工况下引入更完整的目标身份与精位姿信息，并叠加由线缆先验构建的障碍地图，使规划器能够在目标选择、轨迹生成与线缆避障三个环节同时受益。评价指标围绕联合链路输出的可执行性展开，包含识别成功率、误抓次数、线缆接触次数、平均规划时间、预装配位姿误差以及抓取成功率。

\begin{table}[htbp]
  \centering
  \caption{不同遮挡工况下感知---规划联合实验结果}
  \label{tab:ch5:per-occlusion}
  \begin{tabular}{|p{2.3cm}|p{1.6cm}|p{1.9cm}|p{1.8cm}|p{1.8cm}|p{2.0cm}|p{2.1cm}|p{1.8cm}|}
    \hline
    方法 & 遮挡工况 & 识别成功率/\% & 误抓次数/次 & 线缆接触次数/次 & 平均规划时间/s & 预装配位姿误差/mm & 抓取成功率/\% \\
    \hline
    Baseline & 轻度遮挡 & 100 & 0 & 2 & 2.15 & 0.4 & 90 \\
    \hline
    Ours & 轻度遮挡 & 100 & 0 & 0 & 2.42 & 0.3 & 100 \\
    \hline
    Baseline & 中度遮挡 & 85 & 1 & 8 & 3.55 & 1.2 & 55 \\
    \hline
    Ours & 中度遮挡 & 95 & 0 & 0 & 3.28 & 0.8 & 95 \\
    \hline
    Baseline & 重度遮挡 & 40 & 2 & 15 & 8.64 & 4.5 & 15 \\
    \hline
    Ours & 重度遮挡 & 90 & 0 & 1 & 4.82 & 1.1 & 90 \\
    \hline
  \end{tabular}
\end{table}

表 \ref{tab:ch5:per-occlusion} 给出了联合实验的主要量化结果。随着遮挡等级由轻度增加到重度，Baseline 的抓取成功率由 90\,\% 下降到 15\,\%，而本文方法仍保持 90\,\% 的成功率水平；在重度遮挡下，误抓次数从 2 次下降为 0 次，线缆接触次数从 15 次降至 1 次，预装配位姿误差从 4.5\,mm降低到 1.1\,mm，平均规划时间由 8.64\,s 收敛至 4.82\,s。上述对比表明，联合系统在观测不确定性增强与可通行域收缩条件下，能够通过完整接口输入降低错误目标选择概率，并通过三维线缆约束抑制与线缆的近距离擦碰与缠绕风险。

进一步结合轨迹安全性与耗时稳定性进行解释。图 \ref{fig:ch5:safety-gap} 给出了本文方法在舱壁与线缆附近的典型规划轨迹形态，轨迹与关键障碍边界保持了更合理的安全间隙，量化上与线缆接触次数的降低趋势一致，说明规划搜索在引入三维线缆障碍约束后能够稳定选择更安全的通行走廊。

\begin{figure}[htbp]
  \centering
  \fbox{\parbox[c][0.22\textheight][c]{0.9\textwidth}{\centering 图5-8\\（空白占位图）}}
  \caption{规划轨迹与舱壁/线缆安全间隙示意图}
  \label{fig:ch5:safety-gap}
\end{figure}

图 \ref{fig:ch5:occ-compare} 从识别与抓取表现角度对遮挡等级变化进行对比：Baseline 在中度与重度遮挡下表现衰减更快，而本文方法的下降幅度相对更小，体现出对目标关键区域可见性下降的鲁棒适应能力。该结论与表 \ref{tab:ch5:per-occlusion} 中识别成功率和抓取成功率的差异一致，说明联合链路的稳定性主要来自目标身份锁定与预抓取位姿初始化的接口质量。

\begin{figure}[htbp]
  \centering
  \fbox{\parbox[c][0.22\textheight][c]{0.9\textwidth}{\centering 图5-9\\（空白占位图）}}
  \caption{不同遮挡程度下识别与抓取表现对比示意图}
  \label{fig:ch5:occ-compare}
\end{figure}

针对规划耗时，图 \ref{fig:ch5:plan-time-box} 以箱线统计方式比较不同方法的规划时间分布。相较于 Baseline 的更宽分布，本文方法的中位数与离散程度更小，说明在统一实验平台与相同任务触发逻辑下，完整感知先验能够减少无效搜索迭代，并稳定搜索过程。图 \ref{fig:ch5:plan-time-dist} 进一步从分布视角验证该趋势，规划耗时在本文方法中表现出更集中的统计特征。

\begin{figure}[htbp]
  \centering
  \fbox{\parbox[c][0.22\textheight][c]{0.9\textwidth}{\centering 图5-7\\（空白占位图）}}
  \caption{不同方法规划耗时箱线图}
  \label{fig:ch5:plan-time-box}
\end{figure}

如图 \ref{fig:ch5:plan-time-dist} 所示，本文方法在规划耗时上的分布更加集中，体现出感知先验的语义约束不仅改善成功率，也降低了搜索过程对噪声与遮挡误差的敏感性。

\begin{figure}[htbp]
  \centering
  \fbox{\parbox[c][0.22\textheight][c]{0.9\textwidth}{\centering 图5-10\\（空白占位图）}}
  \caption{规划耗时分布统计图}
  \label{fig:ch5:plan-time-dist}
\end{figure}

最后，从失效模式角度解释联合系统为何在重遮挡工况下仍能维持可执行性。图 \ref{fig:ch5:fault-analysis} 给出了典型失败案例：当精位姿估计与线缆三维障碍约束缺失或不一致时，系统更容易出现目标误锁定、轨迹擦碰以及线缆缠绕等问题。该失效形态与表 \ref{tab:ch5:per-occlusion} 中误抓次数与线缆接触次数的变化规律相互印证，说明线缆先验的三维化与目标位姿的接口质量能够共同作用于“选对目标—走对路径—避开线缆”的可执行性链路。

\begin{figure}[htbp]
  \centering
  \fbox{\parbox[c][0.22\textheight][c]{0.9\textwidth}{\centering 图5-11\\（空白占位图）}}
  \caption{典型失败案例分析图}
  \label{fig:ch5:fault-analysis}
\end{figure}

综上，5.3 节实验表明，在狭窄机舱与柔性线缆干扰条件下，通过完整的感知输出—三维障碍映射—规划生成接口输入，联合系统能够在遮挡等级提升时保持更高的抓取成功率，并在重度遮挡下显著降低误抓与线缆接触风险。该结论为后续装配过程的稳定预装配位姿初始化提供了直接实验依据。

\section{面向航空制造的全流程综合验证}
\subsection{实验流程设计}
为验证前述子系统方案在完整系统尺度上的可用性，本节开展端到端综合实验，重点检验：在狭窄机舱几何约束与柔性线缆干扰并存条件下，感知输出是否能够稳定驱动抓取规划；规划轨迹是否能够在时变线缆约束下安全执行；装配控制是否能够在接触非线性出现时维持力响应稳定性，并完成连续任务链闭环。

实验以缩比机舱的多实例混叠为输入场景，同时叠加狭窄通道约束与柔性拖缆干扰。四个连续阶段按动作链组织并形成可重复的接口边界：第一阶段由感知模块完成目标身份锁定、精位姿估计与线缆状态量输出；第二阶段由抓取与规划模块完成接近、夹持、避障搬运以及预对准，并将线缆三维约束纳入可行域评估；第三阶段由装配控制模块完成接触调节、搜孔与插接动作，在力觉闭环下实现误差收敛；第四阶段由状态管理逻辑触发复位与下一目标实例处理，从而考察连续作业能力与阶段切换的稳定性。

\begin{figure}[htbp]
  \centering
  \fbox{\parbox[c][0.22\textheight][c]{0.9\textwidth}{\centering 图5-11\\（空白占位图）}}
  \caption{全流程综合验证实验流程}
  \label{fig:ch5-flow}
\end{figure}

在评价指标选择上，本节以作业节拍与装配力响应一致性为核心量化指标，并以线缆干扰下的成功率检验工程可用性。人工装配用于提供现实工艺条件下的参考上界，用于解释自主系统在效率与稳定性之间的综合权衡。

\subsection{综合效能评估与结果分析}
全流程综合验证在缩比机舱平台上按图 \ref{fig:ch5-flow} 所示的动作链组织执行。每次任务以相同的初始状态触发系统：先完成目标身份锁定与预对准初始化，再进入抓取搬运与基于线缆三维约束的安全轨迹生成，随后进入插接接触阶段并由力觉闭环驱动阻抗调节，最后通过状态管理切换至下一目标实例处理流程。实验过程中记录每次作业节拍、接触阶段的力响应波动形态，并统计线缆干扰条件下的任务完成情况。基于上述采集数据，表 \ref{tab:ch5-fullstat} 汇总了人工装配（熟练工）与本文自主系统在典型作业工况下的统计结果。
\begin{table}[htbp]
  \centering
  \caption{全流程综合性能统计结果}
  \label{tab:ch5-fullstat}
  \begin{tabular}{|p{3.3cm}|p{3.8cm}|p{3.3cm}|}
    \hline
    评估指标 & 人工装配（熟练工） & 本文自主系统 \\
    \hline
    单次作业节拍 / s & 15.5 & 22.8 \\
    \hline
    装配力峰值方差 / N & 2.5 & 0.4 \\
    \hline
    线缆干扰下成功率 / \% & 98 & 96.7 \\
    \hline
  \end{tabular}
\end{table}

从统计指标可见，本文自主系统在节拍上存在一定成本（22.8\,s 相比人工的 15.5\,s），该成本主要来源于全流程链路中对安全间隙与接触稳定性的冗余校验与力觉闭环调节时间。更关键的是，装配力峰值方差由人工装配的 2.5\,N 降至 0.4\,N，表明在插接与搜孔接触非线性阶段，自主系统能够将力响应波动限制在更小范围内，从而减少由接触力突变引发的对准漂移与二次调整次数。在线缆干扰下的插接成功率达到 96.7\%，与人工装配的 98\% 保持接近，说明本文方案能够在刚柔耦合场景中维持可执行的连续动作链。

\par
从时间尺度匹配角度看，全流程节拍差异并不意味着插接过程的误差积累，而更偏向体现在“插接前的安全准备阶段”。在节点触发机制约束下，视觉侧以约 15\,Hz 刷新目标位姿与线缆状态量；规划侧在任务触发条件下生成阶段轨迹并交付执行；控制执行侧以 500\,Hz 输出速度指令，在单次视觉更新间隔内形成多轮闭环校正；进入插接/接触阶段后，力/力矩传感器以 1000\,Hz 采样接触载荷并驱动力觉闭环调节。该多速率协同使得“感知更新—轨迹约束—力觉校正”在同一工况窗口内保持稳定对应关系，从而抑制非线性接触引起的末端轨迹抖动。

\begin{figure}[htbp]
  \centering
  \fbox{\parbox[c][0.22\textheight][c]{0.9\textwidth}{\centering 图5-12\\（空白占位图）}}
  \caption{全流程自主装配实验关键帧序列图}
  \label{fig:ch5-keyframes}
\end{figure}

如图 \ref{fig:ch5-keyframes} 所示，全流程自主装配的关键帧序列可划分为预对准、搬运避障与插接接触三个动作段。预对准阶段的末端姿态连续性由上游抓取规划输出保证，使接近过程能够在不破坏线缆空间约束的前提下完成姿态导入；搬运与避障阶段通过将线缆状态量转化为规划可行域约束，降低末端轨迹对线缆附近几何边界的扰动概率；插接阶段进入接触非线性后，力觉反馈驱动控制律对接触误差进行实时校正，并在阻抗调节框架下抑制力波动的传播。上述三个阶段的动作边界与时序衔接方式与图 \ref{fig:ch5-flow} 所示流程设计一致，因此关键帧序列形态与力峰值方差统计能够相互印证。

\par
接触阶段的力响应稳定性可从阻抗调节的等效动力学关系得到解释。末端在约束环境下的期望位移误差可由阻抗模型对应到外载荷响应，其典型形式为

$$M\ddot{x} + B\dot{x} + Kx = f_{ext}$$

其中，$x$ 为末端等效位移（相对插接参考路径），$M$ 表示等效质量，$B$ 表示等效阻尼，$K$ 表示等效刚度，$f_{ext}$ 为外部接触力/力矩通过坐标映射得到的等效外载荷。该模型将插接过程中外载荷引起的扰动映射为阻抗系统的动态响应，从而在控制过程中抑制接触非线性导致的力峰值突变，进而降低装配力峰值方差并提升收敛性。

如图 \ref{fig:ch5-fig19} 所示，极端打结工况会同时改变线缆受力方向并重构局部几何约束，使原有线缆方向先验与可通行域引导的可靠性降低，规划阶段需要更多依赖接触阶段的力觉闭环来实现误差收敛。图 \ref{fig:ch5-fig18} 的综合效能结果表明，即便线缆约束一致性下降，自主系统仍保持较高成功率与力响应稳定性，其原因在于：接触阶段的力觉闭环能够对接触调节误差进行在线修正，而阻抗调节机制能够限制力波动在末端与接触界面的扩散范围，从而避免由几何先验失配引发的插接失稳。

\par
在极端打结场景下，主要失效风险来源于两类误差耦合：其一，线缆三维约束的局部形态变化会导致规划可行域边界的估计偏移；其二，线缆受力方向变化会改变接触界面的摩擦与支撑条件，使得插接几何约束更易出现非理想偏转。力觉闭环通过对接触误差进行在线修正，把上述偏移从“开环几何导向”转化为“闭环动力响应”，因此仍能够在不稳定先验条件下维持插接过程的可收敛性。

\begin{figure}[htbp]
  \centering
  \fbox{\parbox[c][0.22\textheight][c]{0.9\textwidth}{\centering 图5-18\\（空白占位图）}}
  \caption{极端打结工况下综合效能对比示意图}
  \label{fig:ch5-fig18}
\end{figure}

\begin{figure}[htbp]
  \centering
  \fbox{\parbox[c][0.22\textheight][c]{0.9\textwidth}{\centering 图5-19\\（空白占位图）}}
  \caption{极端打结工况下线缆受力方向与局部几何约束变化示意图}
  \label{fig:ch5-fig19}
\end{figure}

为进一步说明全流程结果的可重复性，本节将前述感知—规划联合验证的定量支撑作为接口质量背景。表 \ref{tab:ch5:per-occlusion} 中，本文在轻度、中度与重度遮挡工况下的抓取成功率分别为 100\%、95\% 与 90\%，并且在重度遮挡下平均规划时间由 2.42\,s 增长至 4.82\,s。该上游接口表现表明，即便观测不确定性增强，系统仍能维持足够稳定的预对准初始化质量，从而使插接阶段的力闭环在更接近目标姿态的起始条件下收敛。因此，全流程实验中的力响应稳定性与高成功率不仅由控制策略保证，也由前端感知与规划链路共同提供了可控的起始误差边界。

\section{本章小结}
本章围绕系统集成、联合验证与全流程评估展开，构建了面向航空制造场景的缩比机舱实验平台，完成手眼标定与 ROS 集成，重投影误差为 0.42 pixels、空间误差约为 0.35 mm，感知链路平均耗时约 65\,ms，可支撑系统闭环运行。感知---规划联合实验表明，目标身份、精位姿与 CableMask 先验能够有效提升抓取锁定、预抓取初始化与线缆避障建图效果；装配实验中，本文方法的峰值接触力降至 $6.5\pm1.2$\,N，成功率达 100\%。全流程综合验证表明，本文系统单次节拍为 22.8\,s，线缆干扰下成功率为 96.7\%，并在狭窄空间与柔性线缆约束条件下表现出较高一致性、自主性与鲁棒性。

